% Civic.ai RAG Thesis - Phase 2 Poster
% Based on the Gemini beamerposter template style.
% Compile with XeLaTeX or LuaLaTeX if Bangla examples are kept.
% Required template files: beamerthemegemini.sty and beamercolorthemecam.sty

\documentclass[final]{beamer}

% ====================
% Packages
% ====================

\usepackage[size=custom,width=120,height=72,scale=1.0]{beamerposter}
\usetheme{gemini}
\usecolortheme{cam}

\usepackage{fontspec}
\setsansfont{Latin Modern Sans}
\setmainfont{Latin Modern Roman}
\newfontfamily\bengalifont{Kohinoor Bangla}

\usepackage{graphicx}
\usepackage{booktabs}
\usepackage{array}
\usepackage{tikz}
\usepackage{pgfplots}
\usepackage{amsmath}
\usepackage{hyperref}
\usepackage{anyfontsize}
\pgfplotsset{compat=1.18}

% ====================
% Lengths
% ====================

\newlength{\sepwidth}
\newlength{\colwidth}
\setlength{\sepwidth}{0.025\paperwidth}
\setlength{\colwidth}{0.30\paperwidth}

\newcommand{\separatorcolumn}{\begin{column}{\sepwidth}\end{column}}
\newcommand{\bn}[1]{{\bengalifont #1}}

% ====================
% Title
% ====================

\title{Civic.ai: A Bangla-English RAG System for Government-Service Question Answering}

\author{Shihab Sarker \and Team Members}

\institute[shortinst]{Department of Computer Science and Engineering \\ Final Year Thesis - Phase 2 Design and Evaluation}

% ====================
% Footer
% ====================

\footercontent{
  \href{https://github.com/Shihabsarker93/civic_ai_rag}{github.com/Shihabsarker93/civic\_ai\_rag} \hfill
  P2: Design Alternatives, Functional Verification, and Preliminary Evaluation \hfill
  Birth and Death Registration Domain}

% ====================
% Body
% ====================

\begin{document}

\begin{frame}[t]
\begin{columns}[t]
\separatorcolumn

% ====================
% Column 1
% ====================

\begin{column}{\colwidth}

  \begin{block}{Problem and Motivation}

    Bangladeshi government-service information is distributed across noisy
    PDFs, HTML pages, notices, legal rules, tables, FAQs, and scanned or
    converted documents. Users ask questions in Bangla, English, and
    code-mixed language, but official information is difficult to search and
    easy to misinterpret.

    \vspace{0.5em}
    \textbf{Thesis goal:} build a factual, explainable, bilingual
    Retrieval-Augmented Generation (RAG) chatbot where answers are grounded in
    retrieved government-service evidence.

    \vspace{0.5em}
    \textbf{Current Phase 2 domain:} Birth and Death Registration.

    \begin{itemize}
      \item Accuracy and source grounding are more important than fluency.
      \item Bangla retrieval is treated as a first-class requirement.
      \item The system must be modular so future domains such as Passport and
        BRTA can be added without creating separate repositories.
    \end{itemize}

  \end{block}

  \begin{alertblock}{Phase 2 Objective Mapping}

    P2 requires multiple engineering/theoretical design solutions and
    functional verification under project constraints.

    \begin{itemize}
      \item \textbf{Design process:} modular RAG pipeline from data cleaning to
        answer generation.
      \item \textbf{Alternative solutions:} Dense-only, BM25-only, Hybrid RRF,
        Hybrid+Rerank, Simple RAG, CivicRAG.
      \item \textbf{Functional verification:} 20-question Bangla evaluation set
        with expected source chunks.
      \item \textbf{Constraints:} local LLMs, Bangla documents, noisy
        government data, no API-key dependency, explainability.
    \end{itemize}

  \end{alertblock}

  \begin{block}{Data and Preprocessing}

    The active domain contains cleaned birth/death registration content from:

    \begin{itemize}
      \item HTML Unicode pages and FAQs.
      \item Converted DOCX/PDF government documents.
      \item Legal acts, rules, notices, fee tables, and application processes.
      \item Manually/semi-manually cleaned text to remove navigation, footers,
        sidebars, repeated news, and irrelevant web noise.
    \end{itemize}

    \vspace{0.4em}
    \textbf{Chunking strategy:} document-type-aware chunking.

    \begin{table}
      \centering
      \small
      \begin{tabular}{p{0.34\linewidth} p{0.52\linewidth}}
        \toprule
        \textbf{Document Type} & \textbf{Chunking Choice} \\
        \midrule
        FAQ JSON & One Q/A pair per chunk \\
        Legal rules/acts & Section or rule-based chunks \\
        Application process & Heading/step-based chunks \\
        Fee table & Table-level and row-level chunks \\
        Notice/OCR content & Heading and semantic block chunks \\
        \bottomrule
      \end{tabular}
    \end{table}

    \textbf{Important design choice:} each chunk separates
    \texttt{retrieval\_text} from original answer evidence.

  \end{block}

\end{column}

\separatorcolumn

% ====================
% Column 2
% ====================

\begin{column}{\colwidth}

  \begin{block}{Proposed CivicRAG Pipeline}

    \begin{center}
    \begin{tikzpicture}[
      node distance=0.62cm,
      every node/.style={font=\small},
      box/.style={rectangle, rounded corners, draw=black!45, fill=blue!7,
        align=center, minimum width=8.7cm, minimum height=0.72cm},
      arrow/.style={->, thick, draw=black!60}
    ]
      \node[box] (raw) {Raw Government Data};
      \node[box, below=of raw] (clean) {Manual / Semi-manual Cleaning};
      \node[box, below=of clean] (chunk) {Document-Type-Aware Chunking + Metadata};
      \node[box, below=of chunk] (dual) {retrieval\_text for Search + Original Content for Grounding};
      \node[box, below=of dual] (embed) {BGE-M3 Embeddings + ChromaDB Vector Index};
      \node[box, below=of embed] (query) {User Query: Bangla / English / Code-mixed};
      \node[box, below=of query] (hybrid) {Dense Search + BM25 Search};
      \node[box, below=of hybrid] (rrf) {RRF Fusion + Reranking / Intent Boosting};
      \node[box, below=of rrf] (safe) {Safe Answer Path or Local LLM Fallback};
      \node[box, below=of safe] (final) {Final Answer + Source IDs + Official Links};

      \draw[arrow] (raw) -- (clean);
      \draw[arrow] (clean) -- (chunk);
      \draw[arrow] (chunk) -- (dual);
      \draw[arrow] (dual) -- (embed);
      \draw[arrow] (embed) -- (query);
      \draw[arrow] (query) -- (hybrid);
      \draw[arrow] (hybrid) -- (rrf);
      \draw[arrow] (rrf) -- (safe);
      \draw[arrow] (safe) -- (final);
    \end{tikzpicture}
    \end{center}

    \vspace{0.4em}
    CivicRAG is not only an LLM wrapper. It includes retrieval design,
    reranking, intent-level safe answer paths, language control, and source
    traceability.

  \end{block}

  \begin{block}{Retrieval Alternatives Compared}

    We compare retrieval methods before finalizing the architecture.

    \begin{enumerate}
      \item \textbf{Dense-only:}
        query embedding $\rightarrow$ ChromaDB cosine similarity.
      \item \textbf{BM25-only:}
        tokenized query $\rightarrow$ sparse lexical matching.
      \item \textbf{Hybrid RRF:}
        combine Dense and BM25 ranked lists using Reciprocal Rank Fusion.
      \item \textbf{Hybrid + Rerank:}
        fused candidates $\rightarrow$ reranking stage $\rightarrow$ top evidence.
    \end{enumerate}

    \vspace{0.4em}
    RRF score:
    \[
      \mathrm{RRF}(d) =
      \sum_{r \in R}
      \frac{w_r}{k + \mathrm{rank}_r(d)}
    \]

    A chunk appearing high in both Dense and BM25 receives a stronger score,
    but a strong chunk from only one retriever can still survive.

  \end{block}

  \begin{exampleblock}{Evaluation Metrics}

    Let $G$ be expected relevant source chunk IDs and $R_k$ be top-$k$
    retrieved chunks.

    \begin{align*}
      \mathrm{Recall@k} &= \mathbf{1}[R_k \cap G \neq \emptyset] \\
      \mathrm{ContextPrecision@k} &= \frac{|R_k \cap G|}{k} \\
      \mathrm{RR} &= \frac{1}{\text{rank of first relevant chunk}} \\
      \mathrm{MRR} &= \frac{1}{N}\sum_{i=1}^{N}\mathrm{RR}_i \\
      \mathrm{nDCG@k} &= \frac{\mathrm{DCG@k}}{\mathrm{IDCG@k}}
    \end{align*}

    Generation metrics are planned using semantic correctness, token F1,
    answer relevance, faithfulness proxy, source citation correctness, and
    manual qualitative scoring.

  \end{exampleblock}

\end{column}

\separatorcolumn

% ====================
% Column 3
% ====================

\begin{column}{\colwidth}

  \begin{block}{P2 Retrieval Results}

    Evaluation set: 20 manually drafted Bangla scenario/paraphrase questions
    with expected source chunk IDs.

    \begin{table}
      \centering
      \small
      \begin{tabular}{lrrrr}
        \toprule
        \textbf{Variant} & \textbf{R@1} & \textbf{R@3} & \textbf{R@5} & \textbf{MRR} \\
        \midrule
        BM25-only & 0.450 & 0.550 & 0.800 & 0.546 \\
        Dense-only & \textbf{0.800} & \textbf{0.900} & \textbf{0.900} & \textbf{0.842} \\
        Hybrid RRF & 0.700 & 0.800 & 0.850 & 0.752 \\
        Hybrid+Rerank & 0.650 & 0.800 & 0.850 & 0.738 \\
        \bottomrule
      \end{tabular}
    \end{table}

    \begin{table}
      \centering
      \small
      \begin{tabular}{lrr}
        \toprule
        \textbf{Variant} & \textbf{Context Precision@5} & \textbf{nDCG@5} \\
        \midrule
        BM25-only & 0.220 & 0.481 \\
        Dense-only & \textbf{0.310} & \textbf{0.731} \\
        Hybrid RRF & 0.300 & 0.661 \\
        Hybrid+Rerank & 0.290 & 0.664 \\
        \bottomrule
      \end{tabular}
    \end{table}

    \textbf{Interpretation:} Dense-only is currently the strongest retrieval
    baseline. This suggests BGE-M3 semantic retrieval is highly effective for
    the current cleaned Bangla domain data. Hybrid/RRF remains a candidate but
    requires tuning.

  \end{block}

  \begin{alertblock}{Key Finding from Paraphrase Testing}

    Same-intent paraphrases must retrieve and answer consistently.

    \vspace{0.3em}
    Example intent: lost birth certificate copy.

    \begin{itemize}
      \item \bn{জন্ম নিবন্ধন হারালে কী করব?}
      \item \bn{আমার জন্মসনদ হারিয়ে গেছে, নতুন কপি পেতে হলে কোথায় যেতে হবে?}
    \end{itemize}

    A failure showed that retrieval could find the correct evidence while
    answer routing selected the wrong procedural answer. We fixed this by
    improving intent detection for lost-certificate paraphrases and adding a
    regression test.

    \vspace{0.3em}
    \textbf{Lesson:} retrieval metrics and answer correctness must be evaluated
    separately.

  \end{alertblock}

  \begin{block}{Design Decisions and Next Steps}

    \textbf{Why not use LLM-only?}
    LLM-only answers risk hallucination and cannot guarantee official source
    grounding.

    \textbf{Why not add all domains immediately?}
    More domains increase ambiguity. Multi-domain Civic.ai needs domain routing:
    \[
      \text{Query} \rightarrow \text{Domain Detection} \rightarrow
      \text{Intent Detection} \rightarrow \text{Domain-specific Retrieval}
    \]

    \textbf{Immediate next steps:}
    \begin{itemize}
      \item Expand the gold evaluation set from 20 to 50--100 reviewed
        questions.
      \item Add paraphrase clusters for each important civic intent.
      \item Tune RRF weights, e.g., dense-heavy settings.
      \item Run full cross-encoder reranking on selected subsets.
      \item Add manual qualitative scoring for faithfulness, correctness,
        completeness, and language quality.
    \end{itemize}

  \end{block}

  \begin{block}{References and Artifacts}

    \small
    \begin{itemize}
      \item Source code and experiment reports:
      \href{https://github.com/Shihabsarker93/civic_ai_rag}{github.com/Shihabsarker93/civic\_ai\_rag}
      \item Evaluation dataset:
      \texttt{domains/birth\_death\_registration/data/evaluation/p2\_20q\_eval\_set.json}
      \item Current report:
      \texttt{docs/evaluation/p2\_experiments/p2\_experiment\_report.md}
      \item Inspired by NextRAG-style comparison: baseline RAG vs proposed RAG.
    \end{itemize}

  \end{block}

\end{column}

\separatorcolumn
\end{columns}
\end{frame}

\end{document}
